\section{Homotopy Type Theory}

Homotopy Type Theory (HoTT) extends MLTT by adopting a homotopical interpretation
of types.
It is enriched by principles such as the \emph{univalence axiom} and \emph{higher inductive types}, 
enabling the direct construction of spaces with specified higher path structure and supporting reasoning up to homotopy.

\subsection{A Homotopical Intepretation}
In HoTT, types can be interpretated either topologically as spaces, 
or as (weak) $\infty$-groupoids from higher-dimensional category theory.
A topological space weakly homotopy equivalent to a CW-complex can be classified upto homotopy type by considering its \textit{fundamental
(weak) $\infty$-groupoid}.~\cite{lurie, quillen2006homotopical} This follows from the fundamental (weak) $\infty$-groupoid construction being adjoint 
to the geometric realization functor. (This is the \textit{Homotopy Hypothesis/Theorem}, introduced by Grothendieck in \emph{Pursuing Stacks}~\cite{grothstacks}).

\subsection{Identity Types}
The precise manner in which types correspond to spaces or $\infty$-groupoids is demonstrated by
establishing an associated $\infty$-groupoid.

In the homotopy interpretation, terms of a type are thought of a points, and the identity type $\mathsf{Id}_A(x,y)$ corresponds to the space of 
paths from $x : A$ to $y : A$ endowed with the compact-open topology. This is exactly the free path space $PA$. $\refl_x$ is the
constant path at $x$.
The iterated path type, $\mathsf{Id}_{\mathsf{Id}_A(x,y)}(p,q)$ corresponds to the iterated free path space $P^2 A$.
Note that $h : \mathsf{Id}_{\mathsf{Id}_A(x,y)}(p,q)$ are homotopies from $p$ to $q$.
And similarly the $n$-th iterated path type corresponds to the $n$-th path space $P^n A$. 

We can extract the (weak) $\infty$-groupoid structure from this observation with a few arguments.

Let $\U$ be fixed and $A : \U$. Define $\Id^0(A) :\deq A$ and $\Id^{\mathsf{succ}(n)} :\deq \Sum{x,y : \Id^n(A)}{(x =y)}$ 
\begin{theorem}
  There is a functorial association of types to weak $\infty$-groupoids, where $A : \U$ maps to the infinity groupoid where $n$-th level objects are terms of 
  $\Id^n(A)$.
\end{theorem}
The idea of the proof is to construct inverse and concatenation operations at a fixed level, and showing they follow
unit, inversion and associativity laws upto terms of a higher level (via path induction). Finally, showing functoriality (across independent
and dependent functions) also follows from path induction.
A full walkthrough for this is laid out in Chapter 2 of \cite{HoTT}.

\subsection{Axioms in HoTT}
Expected semantic behaviour does not directly follow from the machinery of HoTT laid out so far. To enforce desired semantic behaviour
we introduce a few `axioms' into our theory, which are merely inhabitants of some type without any dedicated inference rules governing them. But first,
we make some definitions.

\begin{definition}[Homotopy of maps] 
  Let $f, g : \Prod{x:A}{B(a)}$, then we say $f$ is homotopic to $g$, denoted
  \[
    (f \sim g) :\deq \Prod{x:A}{(f(x) =_{B(x)} g(x))}
  \]
\end{definition}

\begin{definition}[Equivalence]
  Let $A, B : \U$. We say $A$ and $B$ are (homotopy) equivalent if there is an inhabitant of 
  \[
    (A \eqv B) :\deq \Sum{f : A \to B}{\Sum{g : B \to A}{(g \circ f \sim \id_A) \times (f \circ g \sim \id_B)}}
  \]
\end{definition}

The notion of equality for functions is too restrictive to capture their expected behaviour.
Classically, we expect two functions to be equal if they take the same values on all the inputs.
To ensure semantic well-behavedness, we must enforce this as an axiom.

\begin{axiom}[Function Extensionality]
   Let $A : \U, B : A \to \U$ and $f, g : \Prod{x:A}{B(a)}$, then we have
   $    (f = g) \eqv (f \sim g)$. 
\end{axiom}
That is, equality of maps is characterized entirely by their values in the codomain.
In particular, if two functions are homotopic, they are equal.

In~\cite{voevodsky2014equivalence}, Voevodsky proposed a variant of the univalence axiom. In its current form, we have
\begin{axiom}[Univalence]
  For $A, B : \U$, there is an equivalence:
  $(A =_{\U} B) \eqv (A \eqv B) $
\end{axiom}
In particular, if $A \eqv B$, then $A =_{\U} B$.
Roughly speaking, this means ``isomorphic objects are equal''. 
In a univalent universe $\U$, isomorphic objects/equivalent types can be treated as equal and one can transport 
(theorems, proofs and properties) across them.

\subsection{Higher Inductive Types}

We can now expand to types generated by constructors along with relations.
We let the constructors output higher-dimensional paths in the type. This way, 
free generation is equivalent to generation on points quotiented by equivalence relations corresponding to these higher paths.

Such a \textit{higher inductive type} primitive in HoTT is the circle type $\Sph^1$ generated by the constructors
\begin{itemize}
  \item A point $\base : \Sph^1$
  \item A path $ \mathsf{loop} : \base =_{\Sph^1} \base $
\end{itemize}
This corresponds to a chosen basepoint $\base$, and a nontrivial path $\mathsf{loop}$ connected the basepoint,
giving a circle.

Many familiar mathematical objects such as the integers, CW-complexes etc. can be built using
this schema. Various approaches exist to formalize such a schema, but have only yeilded partial satisfaction. 
Some of the existing approaches are using Quotient Inductive Types~\cite{altenkirch2016type}, Homotopy Initial Algebras~\cite{sojakova2015higher} , and computationally
most satisfactory, Cubical Type Theory~\cite{cohen2016cubical, cavallo2021higher}. A completely satisfactory formalism remains an open problem.
